% ============================================================================
% PREÁMBULO - Configuraciones y Paquetes (Alineado a Guía UNAP 2.0 y APA 7)
% ============================================================================

% Clase de documento
\documentclass[12pt,oneside]{book}

% -------------------------------------------------------------------------
% Paquetes fundamentales y Tipografía
% -------------------------------------------------------------------------

% Idioma español - Traducir "Cuadro" como "Tabla"
\usepackage[spanish,es-tabla]{babel}
\addto\captionsspanish{\def\tablename{Tabla}}

% Codificación UTF-8
\usepackage[utf8]{inputenc}

% Paquete para geometría de páginas (Márgenes UNAP)
\usepackage{geometry}

% Paquete para sangría francesa en bibliografía APA 7
\usepackage{hanging}

% Paquete para ecuaciones matemáticas básicas
\usepackage{amsmath}
\usepackage{amssymb}
\usepackage{amsfonts}
\usepackage{bm}

% Paquetes para diagramas, esquemas y tablas anchas
\usepackage{tikz}
\usetikzlibrary{shapes.geometric, arrows, positioning, fit, calc, backgrounds, babel}
\usepackage{pgfplots}
\pgfplotsset{compat=1.18}
\usepackage{pdflscape}

% Fuente Times New Roman (12pt según Guía UNAP)
\usepackage{mathptmx}

% Interlineado 1.5 (Configurado según estándar UNAP)
\usepackage{setspace}
\setstretch{1.5}

% Márgenes Oficiales UNAP:
% Izquierdo: 3.0 cm (para empastado)
% Derecho: 2.5 cm | Superior: 2.5 cm | Inferior: 2.5 cm
\geometry{
    left=3.0cm,
    right=2.5cm,
    top=2.5cm,
    bottom=2.5cm
}

% -------------------------------------------------------------------------
% Control de Párrafos y Formato
% -------------------------------------------------------------------------

% Control de líneas huérfanas/viudas
\widowpenalty=10000
\clubpenalty=10000

% Sangría de primera línea (1.25 cm según APA 7 / Guía UNAP)
\usepackage{indentfirst}
\setlength{\parindent}{1.25cm}

% Saltos de página controlados
\usepackage{emptypage}

% Formato de encabezados y pies de página (Numeración en pie de página centrado)
\usepackage{fancyhdr}
\fancyhf{}
\fancyhead[C]{}
\fancyfoot[C]{\thepage}
\pagestyle{fancy}
\renewcommand{\headrulewidth}{0pt}

% Hipervínculos en el PDF
\usepackage{url}
\usepackage[colorlinks=true,linkcolor=black,anchorcolor=black,citecolor=black,filecolor=black,urlcolor=black]{hyperref}

% -------------------------------------------------------------------------
% Configuración de Tablas y Figuras (Guía UNAP 2.0 / APA 7)
% -------------------------------------------------------------------------

\usepackage{graphicx}
\graphicspath{{figuras/}}

\usepackage{float}
\usepackage{array}
\usepackage{tabularx}
\usepackage{booktabs}
\usepackage{caption}

% Estilo de Leyendas (Captions) alineado a Guía UNAP págs. 21 y 22:
% "Tabla 1" en Negrita, Salto de línea, y Título descriptivo en Cursiva
\DeclareCaptionFormat{unap}{\textbf{#1}\\\textit{#3}}
\captionsetup{
    format=unap,
    justification=raggedright,
    singlelinecheck=false,
    labelfont=bf,
    font=normalsize,
    skip=5pt
}

% -------------------------------------------------------------------------
% Configuración de Títulos y Capítulos (Guía UNAP 2.0)
% -------------------------------------------------------------------------

\usepackage{titlesec}

% CAPÍTULO I (16pt, Negrita, Centrado)
% TÍTULO DEL CAPÍTULO (14pt, Negrita, Centrado, MAYÚSCULA)
\titleformat{\chapter}[display]
    {\bfseries\centering}
    {\Large\MakeUppercase{CAPÍTULO \Roman{chapter}}}
    {5pt}
    {\large\MakeUppercase}
    []

\titlespacing*{\chapter}{0pt}{-10pt}{20pt}

% NIVEL 2 (Ej. 1.1. TÍTULO SECCIÓN) - MAYÚSCULA, Negrita, Izquierda
\titleformat{\section}
    {\bfseries\normalsize\raggedright}
    {\thesection}
    {0.5em}
    {\MakeUppercase}

% Configuración de profundidad de numeración y tabla de contenidos (Nivel 4: 2.2.1.1)
\setcounter{secnumdepth}{3}
\setcounter{tocdepth}{3}

% NIVEL 3 (Ej. 2.2.1. Título Subsección) - Minúscula, Negrita, Izquierda
\titleformat{\subsection}
    {\bfseries\normalsize\raggedright}
    {\thesubsection}
    {0.5em}
    {}

% NIVEL 4 (Ej. 2.2.1.1. Título Sub-subsección) - Minúscula, Negrita, Izquierda
\titleformat{\subsubsection}
    {\bfseries\normalsize\raggedright}
    {\thesubsubsection}
    {0.5em}
    {}

% -------------------------------------------------------------------------
% Listas y Viñetas
% -------------------------------------------------------------------------
\usepackage{enumitem}
\setlist[itemize]{leftmargin=1.25cm, topsep=2pt, itemsep=2pt}
\setlist[enumerate]{leftmargin=1.25cm, topsep=2pt, itemsep=2pt}